\section{Botnets}

A \textbf{botnet} is a collection of compromised machines, called
\emph{bots} or \emph{zombies}, controlled remotely by an attacker called
the \emph{botmaster}. Machines can be recruited through malware, worms,
phishing emails, pirated software, or attacks against weak passwords.

\paragraph{Command and Control}
Once infected, bots communicate with one or more
\emph{Command-and-Control} (C\&C) servers. These servers allow the
botmaster to send commands, install updates, or collect stolen data such
as passwords, credit card numbers, or private keys.

\paragraph{Uses of Botnets}
Botnets can be used to send spam, perform DDoS attacks, steal
credentials using keyloggers, distribute malware or ransomware, perform
crypto mining, or automate fraudulent actions.

\paragraph{Architectures}
Botnets can be centralized, decentralized, or hybrid. In a centralized
architecture, bots receive commands from a small number of C\&C servers.
In a decentralized architecture, bots communicate in a peer-to-peer way,
which makes the botnet harder to take down.

\paragraph{Resilience Techniques}
Modern botnets try to hide their infrastructure. For example, they may
use \emph{Fast Flux DNS}, where the IP address of a domain changes
frequently, or domain generation algorithms, where bots try many possible
domain names to find an active C\&C server.

\paragraph{Fighting Botnets}
Defending against botnets is difficult. Possible strategies include
preventing infection, taking down C\&C servers, seizing domain names, or
detecting suspicious behavior such as many failed DNS queries. However,
botnets are often spread across several countries, which makes legal and
technical takedown operations complex.

\paragraph{Attack Economy}
Botnets are also part of a criminal market. Attackers can rent botnets,
buy stolen accounts or credit card data, or purchase DDoS attacks as a
service, often called \emph{booters} or \emph{stressers}.